% Переработанный раздел для noether_casimir_drive.tex
% Заменяет блок от "ON THE INCOMPLETENESS OF THE STANDARD MOMENTUM BALANCE" до "CONCLUSION".
% Включает: замечание к сноске 22 (Tamm_ponderomotive_forces_update.tex, §33) и
% Appendix B статьи Casimir_honeycomb_drive.tex.

\vspace{\myvspacebeforesubsection}
\subsection*{\centering{\fontencoding{T1}\fontfamily{ptm}{\selectfont ON THE INCOMPLETENESS OF THE STANDARD MOMENTUM BALANCE}}}
\vspace{\myvspaceaftersubsection}

It is essential to specify \emph{which} momentum is not conserved, so that the
result is not misread as a violation of a fundamental conservation law. The
standard momentum-balance equation of continuum electrodynamics (Tamm,
\S\,105) reads
\end{multicols}
\[
\nabla\cdot\mathbf{T} = \frac{\partial\mathbf{g}}{\partial t} + \mathbf{f}_{\mathrm{vol}},
\qquad
\mathbf{g} = \frac{1}{c^2}\mathbf{S} = \frac{1}{4\pi c}[\mathbf{E}\times\mathbf{H}],
\]
\begin{multicols}{2}
where $\mathbf{T}$ is the Maxwell stress tensor, $\mathbf{g}$ the field
momentum density and $\mathbf{f}_{\mathrm{vol}}$ the volume (ponderomotive)
force density. For the vacuum fluctuations of a static honeycomb system the
mean Poynting vector vanishes, $\langle\mathbf{S}\rangle = 0$, hence
$\langle\partial\mathbf{g}/\partial t\rangle = 0$.

\textbf{Where the standard derivation is incomplete.}
The passage from volume forces to stresses rests on the Gauss theorem, which
Tamm applies (\S\,33, eq.~33.7) under the explicit condition that the
components of the stress tensor are \emph{continuous} inside the volume; surfaces
of discontinuity are excluded there, and only their surface \emph{charges} are
treated later (\S\,34, eq.~34.8). A vacuum--metal boundary of a resonator can carry
zero surface charge and still produce a finite jump of the stress tensor: the tensor
evaluated with the fields on the vacuum side differs from the one evaluated with the
fields inside the wall. On the two sides of such a surface element the stresses are
then \emph{not} equal and opposite. This case is not considered in \S\,33, nor
anywhere in Tamm's book, including \S\,105. For a volume $V$ that contains internal
discontinuity surfaces $S_j$ the Gauss theorem has to be written as
\end{multicols}
\[
\int_V \frac{\partial T_{ik}}{\partial x_k}\,dV
= \oint_{\partial V_{\mathrm{ext}}} T_{ik}\,n_k\,dS
\;-\; \sum_j \int_{S_j} [\![\,T_{ik}\,]\!]\,n_k\,dS_j ,
\]
\begin{multicols}{2}
where $[\![\cdot]\!]$ is the jump across $S_j$. The customary identification of the
field momentum with $\mathbf{g}=\mathbf{S}/c^2$ keeps only the external surface
integral and drops the jump terms.

\textbf{The jump on a honeycomb plate.}
The plate has two such internal boundaries: the bottoms of the cells on the
structured side and the smooth face on the other side. The stress on each of them is
the electromagnetic pressure of the vacuum modes. Appendix~B of \cite{Drozdov2024}
computes this pressure on the $xy$ face of an $a\times b\times h$ rectangular
resonator, $f_z/S=\int(H_x^2+H_y^2)\,dS/(8\pi S)$, mode by mode, and relates it
to the mode energy density $E/V$. For a wave propagating along $z$
\end{multicols}
\[
\frac{f_z/S}{E/V}=1 \quad\text{(electric mode, } H_z=0\text{)},
\qquad
\frac{f_z/S}{E/V}=\frac{a^2k_z^2}{a^2k_z^2+\pi^2(n_x^2+n_y^2)} \quad\text{(magnetic mode, } E_z=0\text{)},
\]
\begin{multicols}{2}
and for every mode type $\lim_{k_z\to\infty}(f_z/S)/(E/V)=1$. Summed over the
zero-point modes, the jump of $\langle T_{zz}\rangle$ across each boundary is
therefore fixed by the vacuum energy density of the adjacent cavity, and the
difference of the two jumps gives
\[
\frac{F}{S}=[\![\langle T_{zz}\rangle]\!]_{\mathrm{honey}}-[\![\langle T_{zz}\rangle]\!]_{\mathrm{smooth}}
\;\ge\;\delta\frac{E}{V}=\hbar c\,\pi\frac{R}{a^4},
\]
the inequality of \cite{Drozdov2024} (Appendix~B). It coincides with the
Hamiltonian result $\partial\langle H\rangle/\partial l$ of the previous section.
For a gold nano-honeycomb with cell size $a\approx 2~\mu$m this gives
$F/S\approx 8.55~\text{dyn/m}^2 = 8.55\times10^{-5}~\text{N/m}^2$.
The \S\,105 reading, completed by the jump terms, therefore does not re-derive the
force but \emph{locates} it: the unbalanced momentum sits in the jump of the vacuum
stress tensor on the internal boundaries, which the reduction to
$\mathbf{g}=\mathbf{S}/c^2$ drops.

\textbf{Equivalent pillbox form.}
Enclosing the whole plate by a pillbox $\partial V_{\mathrm{vac}}$ with faces in the
free vacuum at $z=\pm d$ (lateral faces cancel by in-plane symmetry) turns the sum of
the jumps into an external flux,
$F_z=\oint_{\partial V_{\mathrm{vac}}}\langle T_{zz}\rangle\,da
=\big[\langle T_{zz}\rangle_{z=+d}-\langle T_{zz}\rangle_{z=-d}\big]S$.
This is the flux of \emph{field} momentum through a purely vacuum boundary, non-zero
solely because the mode structure, and hence $\langle T_{zz}\rangle$, differs on the
honeycomb and on the smooth side.

Two remarks make the interpretation unambiguous.
First, it is the momentum of the \emph{field} (equivalently, of the material plate that
absorbs it), not the total momentum of the closed laboratory system, that is not
conserved; the total momentum of plate, field and support is conserved, exactly as for
an ordinary sail in a wind. Second, the renormalised vacuum stress
$\langle T_{\mu\nu}\rangle$ is a physically measured quantity --- the Casimir effect
itself --- so the jump terms require no assumption of a material aether: their value
is fixed by the same $\langle T_{\mu\nu}\rangle$ that is confirmed experimentally for
the symmetric configuration. The broken translational symmetry of the boundary
conditions, not of space itself, is what renders $\partial\langle H\rangle/\partial l\neq 0$.

